\chapter*{Abstract}
\addcontentsline{toc}{chapter}{Abstract}

\begin{center}
  {\Large\bfseries Intelligent Information Assistant}\\[0.3em]
  {\large\itshape for Users of an Insurance Company}
\end{center}
\vspace{0.5em}

\begingroup
\setlength{\parskip}{0.8em}

Customers of an insurance company often ask the same questions about coverage, exclusions,
deductibles and claims. The answers are in the company's public documents, but these are long
and written in contractual language. The project builds a RAG (Retrieval-Augmented Generation)
assistant that answers at any time of day using only these approved sources, always citing the
one it used.

The goal is to provide consistent information even outside business hours and to reduce the
repetitive requests handled manually.

Starting from the company's official, publicly accessible documentation (FAQs, forms, customer
support pages, anonymized past requests, etc.), the system removes duplicates and extracts the
text.

For each request it produces an answer together with references to the sources used,
a classification of the question and an indication of the confidence level.

If the sources are not sufficient, the assistant does not answer and says so. It classifies
every request, estimates its confidence and proposes a handover to a human operator when human
judgment is needed: complaints, disputes, the assessment of a claim. Every answer is logged with
its sources, confidence level and model version, and the questions that are not covered end up
in an anonymized log.

The assistant does not decide whether a claim is covered or reimbursable: it provides
information and leaves the decision to the operator. The operator receives the request already
summarized, with the cited clauses and the reason for the handover. The operator's corrections
are used to improve the answers and the content.
\par
\endgroup
